% preamble_common.tex
% Gemeinsame Präambel für alle Druck-Editionen (main.tex, main_kdp_print.tex,
% main_bod_print.tex). Enthaelt alles, was NICHT von der Trimgroesse abhaengt.
% \documentclass und \geometry bleiben in der jeweiligen Wrapper-Datei, da sie
% pro Edition unterschiedlich sind.

\usepackage[ngerman]{babel}
\usepackage[T1]{fontenc}
\usepackage{textcomp}
\usepackage[utf8]{inputenc}
\usepackage{lmodern}
\usepackage{microtype}
\usepackage{geometry}

\usepackage{graphicx}
\usepackage{adjustbox}
\usepackage{booktabs}
\usepackage{longtable}
\usepackage{array}
\usepackage{xcolor}
\usepackage{hyperref}
\usepackage{listings}
\usepackage{enumitem}
\usepackage{csquotes}
\usepackage{caption}
\usepackage{float}
\usepackage{tikz}
\usetikzlibrary{positioning,arrows.meta,shadows,calc,quotes,backgrounds}
\usepackage{tcolorbox}
\usepackage{amsmath}
\usepackage{makeidx}
\usepackage[acronym,toc]{glossaries}

\makeindex
\makeglossaries

\hypersetup{
    colorlinks=true,
    linkcolor=blue!50!black,
    urlcolor=blue!50!black,
    citecolor=blue!50!black,
    pdftitle={SysML v2, MBSE und ROS2 für Einsteiger},
    pdfauthor={Wolfgang Becker}
}

\RedeclareSectionCommand[tocnumwidth=3.8em]{section}
\RedeclareSectionCommand[tocnumwidth=4.8em]{subsection}
\RedeclareSectionCommand[tocnumwidth=5.8em]{subsubsection}

\lstdefinestyle{code}{
    basicstyle=\ttfamily\small,
    breaklines=true,
    frame=single,
    numbers=left,
    numberstyle=\tiny,
    columns=fullflexible,
    showstringspaces=false,
    tabsize=2,
    extendedchars=true,
    inputencoding=utf8,
    upquote=true,
    literate=
        {ä}{{\"a}}1 {ö}{{\"o}}1 {ü}{{\"u}}1
        {Ä}{{\"A}}1 {Ö}{{\"O}}1 {Ü}{{\"U}}1
        {ß}{{\ss{}}}1
}
\lstset{style=code}

\newtcolorbox{lernziele}{
    colback=blue!3!white,
    colframe=blue!50!black,
    title=Lernziele
}
\newtcolorbox{praxisbox}{
    colback=green!3!white,
    colframe=green!40!black,
    title=Praxisbeispiel
}
\newtcolorbox{hinweisbox}{
    colback=yellow!8!white,
    colframe=yellow!50!black,
    title=Hinweis
}
\newtcolorbox{uebungbox}{
    colback=gray!5!white,
    colframe=gray!60!black,
    title=Übung
}

\newacronym{mbse}{MBSE}{Model-Based Systems Engineering}
\newacronym{sysml}{SysML}{Systems Modeling Language}
\newacronym{ros2}{ROS2}{Robot Operating System 2}
\newacronym{bdd}{BDD}{Block Definition Diagram}
\newacronym{ibd}{IBD}{Internal Block Diagram}
\newacronym{slam}{SLAM}{Simultaneous Localization and Mapping}
\newacronym{ml}{ML}{Machine Learning}
\newacronym{api}{API}{Application Programming Interface}

\newglossaryentry{requirement}{
    name=Anforderung,
    description={Beschreibung einer Fähigkeit, Eigenschaft oder Einschränkung, die ein System erfüllen muss}
}
\newglossaryentry{block}{
    name=Block,
    description={Zentrales Strukturelement in SysML zur Beschreibung von Komponenten, Subsystemen oder logischen Einheiten}
}
\newglossaryentry{node}{
    name=ROS2-Node,
    description={Eigenständiger Prozess in ROS2, der Daten verarbeitet und über Topics, Services oder Actions kommuniziert}
}
\newglossaryentry{topic}{
    name=Topic,
    description={Nachrichtenkanal in ROS2 für publish-subscribe-Kommunikation}
}
\newglossaryentry{traceability}{
    name=Traceability,
    description={Nachvollziehbare Verbindung zwischen Anforderungen, Architektur, Implementierung und Tests}
}


\title{SysML v2, MBSE und ROS2 für Einsteiger\\
\large Ein praxisnahes Lehrbuch am Beispiel eines autonomen Indoor-Roboters}
\author{Wolfgang Becker}
\date{\today}
